\section{Forschungsfrage}\label{chapter:Forschungsfrage}
\subsection{Dynamische Zugriffskontrolle und Skalierbarkeit von Dokumentensystemen}\label{subsec:dynamische_zugriffskontrolle_und_skalierbarkeit}
Im architektonischen Teil dieser Arbeit steht die Frage im Mittelpunkt, wie ein Dokumentensystem so entworfen werden kann, dass es sowohl sicher als auch langfristig erweiterbar bleibt. Dabei geht es insbesondere um die Verbindung von Zugriffskontrolle, Administrierbarkeit und technischer Skalierbarkeit in einer webbasierten Umgebung.

Aus diesem Themenfeld ergeben sich folgende Forschungsfragen:
\begin{enumerate}
    \item Wie kann ein flexibles Rollen- und Berechtigungssystem entworfen werden, das sich dynamisch an unterschiedliche Benutzergruppen anpasst und dennoch leicht administrierbar bleibt?
    \item Wie lässt sich eine skalierbare und plattformunabhängige Architektur für eine webbasierte Dokumentenablage entwerfen, die große Datenmengen performant verarbeitet und gleichzeitig hohe Sicherheitsanforderungen erfüllt?
\end{enumerate}

Mit diesen Fragen soll analysiert werden, welche Architektur- und Sicherheitsprinzipien für ein modernes Dokumentensystem geeignet sind. Im Fokus stehen dabei die strukturierte Vergabe von Berechtigungen, die zentrale Verwaltung von Identitäten und Zugriffsrechten sowie der Aufbau einer technischen Architektur, die Lastspitzen, wachsende Dokumentbestände und unterschiedliche Zielplattformen zuverlässig bewältigen kann.

\subsection{Automatisierte Dokumentenanalyse und Suchoptimierung}\label{subsec:automatisierte_dokumentenanalyse_und_suchoptimierung}
Im zweiten Schwerpunkt dieser Arbeit wird untersucht, wie Dokumente automatisiert inhaltlich verarbeitet und anschließend effizient wieder aufgefunden werden können. Dabei stehen sowohl die KI-gestützte Erkennung des Dokumenttyps als auch die Ausgestaltung einer geeigneten Such- und Filterarchitektur im Vordergrund.

Aus diesem Themenfeld ergeben sich folgende Forschungsfragen:
\begin{enumerate}
    \item Wie genau können KI-Modelle anhand des Inhalts erkennen, um welchen Dokumententyp es sich handelt?
    \item Welche Architektur für Suche und Filterung bietet die höchste Geschwindigkeit und Genauigkeit bei der Auffindbarkeit von Dokumenten?
\end{enumerate}

Mit diesen Fragen soll erklärt werden, auf welcher technischen Grundlage Dokumente automatisiert klassifiziert und für nachgelagerte Prozesse aufbereitet werden können. Zusätzlich wird analysiert, welche Suchlogiken und Metadatenstrukturen erforderlich sind, damit Dokumente in einer wachsenden Ablage schnell, präzise und benutzerorientiert gefunden und eingegrenzt werden können.
